\chapter{Implementação e Avaliação Experimental}

Este capítulo apresenta a implementação do sistema desenvolvido e os experimentos conduzidos com o objetivo de avaliar seu desempenho, escalabilidade e corretude geométrica. Em particular, busca-se analisar o impacto de diferentes estratégias nas etapas de filtragem e refinamento, bem como comparar os algoritmos SAT e GJK/EPA em um ambiente de execução web, no qual restrições de desempenho e responsividade são relevantes.

Os experimentos foram realizados em um computador com as seguintes especificações:
\begin{itemize}
	\item Processador: AMD Ryzen 5 5500U
	\item Placa de Vídeo: AMD Radeon Vega 7 (integrada)
	\item Memória RAM: 24 GB DDR4
	\item Sistema Operacional: Ubuntu 24.04 LTS
\end{itemize}

\section{Arquitetura e Implementação}

Foi desenvolvida uma biblioteca em \textit{TypeScript} responsável pela implementação dos algoritmos de detecção de colisão, sendo reutilizada tanto na aplicação de visualização quanto no ambiente de testes. A engine permite a execução em dois modos distintos: utilizando o algoritmo do Teorema do Eixo Separador (SAT) ou a combinação dos algoritmos GJK/EPA.

A arquitetura do sistema foi projetada de forma modular, separando a lógica geométrica, os algoritmos de detecção e resposta à colisão e a camada de visualização, delegando esta última à aplicação cliente. Essa separação favorece a reutilização de código e possibilita a comparação direta entre diferentes abordagens de detecção de colisão.

Os objetos convexos são representados como listas de vértices ordenados, associadas a um conjunto de restrições. Cada restrição atua sobre um par de partículas, aplicando um processo de relaxamento que permite modelar tanto estruturas rígidas quanto deformáveis. Essa modelagem possibilita a integração da detecção de colisão com simulações físicas baseadas em partículas. 

Este trabalho disponibiliza uma implementação pública com o objetivo de permitir a reprodução dos experimentos e a extensão das técnicas apresentadas~\footnote{Disponível em: \url{https://github.com/devdiegomatos/simplified-physics-animation-web}}.

% O código-fonte da biblioteca desenvolvida, bem como as aplicações de demonstração, encontram-se disponíveis em: \url{}.
\section{Ambiente de Visualização}

Com o objetivo de permitir a análise qualitativa do comportamento dos algoritmos, foram desenvolvidas duas aplicações interativas (\textit{demos})~\footnote{https://devdiegomatos.github.io/simplified-physics-animation-web/}. utilizando \textit{Vue.js}, \textit{Vite} e \textit{p5.js}

A primeira demonstração consiste em um ambiente contendo múltiplos objetos convexos, com número de vértices variando entre 3 e 8, inseridos em uma caixa com obstáculos. Esse cenário permite a observação da dinâmica de colisões em tempo real. A interface possibilita o controle de parâmetros relevantes, como a quantidade de objetos, o algoritmo de detecção de colisão (SAT ou GJK/EPA) e a estratégia de filtragem, incluindo o uso de uma grade uniforme. A Figura~\ref{fig:demo1} ilustra esse ambiente.

\begin{figure}[htb]
    \centering
    \includegraphics[width=\textwidth]{demo1.png}
    \caption{Ambiente de simulação com múltiplos objetos convexos em uma caixa, permitindo a configuração de parâmetros de colisão e filtragem. Fonte: Autor.}
    \label{fig:demo1}
\end{figure}

A segunda demonstração apresenta simulações de tecidos e treliças, evidenciando a interação entre objetos deformáveis. Nesse cenário, as restrições entre partículas são responsáveis pela manutenção estrutural, enquanto o sistema de colisão garante a interação física entre os elementos. A Figura~\ref{fig:demo2} apresenta um exemplo desse ambiente.

\begin{figure}[htb]
    \centering
    \includegraphics[width=\textwidth]{demo2.png}
    \caption{Simulação de estruturas deformáveis, como tecidos e treliças, utilizando restrições entre partículas e detecção de colisão. Fonte: Autor.}
    \label{fig:demo2}
\end{figure}

\section{Metodologia de Testes}

Os experimentos foram conduzidos com três objetivos principais:
\begin{itemize}
    \item Avaliar o impacto da utilização de uma estrutura de filtragem no desempenho do sistema;
    \item Comparar o desempenho dos algoritmos SAT e GJK/EPA na fase de refinamento;
    \item Analisar a escalabilidade das abordagens em função do número de objetos.
    \item Comparar o desempenho da implementação desenvolvida com a biblioteca \textit{Planck.js},
\end{itemize}
Os testes foram realizados sem a etapa de visualização, a fim de isolar o custo computacional da detecção de colisão. Para garantir a reprodutibilidade dos experimentos, foi utilizada geração pseudoaleatória com semente fixa por meio da biblioteca \textit{seedrandom}, uma vez que a função \textit{Math.random} não permite controle de semente.

O cenário de teste corresponde ao ambiente apresentado na Figura~\ref{fig:demo1}, no qual objetos convexos são posicionados aleatoriamente dentro de uma região delimitada. Foram variados o número de objetos e o tamanho do espaço, de modo a manter aproximadamente constante a densidade de ocupação. Essa estratégia permite avaliar o impacto do crescimento do sistema sem introduzir vieses relacionados à distribuição espacial.

Foram considerados quatro casos experimentais:
\begin{itemize}
    \item SAT sem fase de filtragem;
    \item SAT com grade uniforme na fase de filtragem;
    \item GJK/EPA sem fase de filtragem;
    \item GJK/EPA com grade uniforme na fase de filtragem.
\end{itemize}
O número de objetos foi variado de 1000 a 10000, em incrementos de 1000. Para cada configuração, foram realizadas 180 execuções independentes, sendo computadas as médias das métricas coletadas.

As métricas analisadas incluem o tempo médio de execução, o número de testes de colisão realizados e o número de colisões detectadas, permitindo avaliar tanto o custo computacional quanto a corretude dos algoritmos. Essa abordagem permite a visualização simultânea de todas as configurações, preservando a comparabilidade entre os métodos.

A biblioteca \textit{Planck.js} é um motor de física 2D para JavaScript, reescrito a partir do Box2D, amplamente utilizado no desenvolvimento de jogos digitais. A \textit{Planck.js} emprega uma estrutura de árvore hierárquica dinâmica de volumes delimitadores (AABB) na fase de filtragem, enquanto a detecção de colisão na fase de refinamento baseia-se em algoritmos como GJK e EPA.

Para garantir a comparabilidade dos resultados, foram utilizados cenários equivalentes aos descritos anteriormente, mantendo-se consistentes os parâmetros de número de objetos, distribuição espacial e número de execuções. Devido a forma como a biblioteca é implementando as foi analisada apenas a métrica de tempo médio de execução.

\section{Resultados Experimentais}

Os resultados obtidos permitem analisar o comportamento dos algoritmos sob diferentes configurações e validar empiricamente as hipóteses estabelecidas.

\begin{figure}[htb]
    \centering
    \includegraphics[width=0.85\textwidth]{comparativo-testes-colisao.png}
    \caption{Número médio de testes de colisão em função do número de objetos para diferentes configurações: SAT e GJK/EPA, com e sem utilização de grade uniforme. Fonte: Autor.}
    \label{fig:testes-comparativo}
\end{figure}

A Figura~\ref{fig:testes-comparativo} evidencia o impacto direto da etapa da filtragem na redução do número de testes de colisão. A utilização da grade uniforme reduz drasticamente a quantidade de pares avaliados, refletindo diretamente na melhoria de desempenho observada. Devido à diferença de ordens de grandeza entre os cenários com e sem fase de filtragem, os resultados são apresentados em escala logarítmica no eixo vertical.

\begin{figure}[htb]
    \centering
    \includegraphics[width=0.85\textwidth]{comparativo-tempo-execução.png}
    \caption{Tempo médio de execução em função do número de objetos para diferentes configurações: SAT e GJK/EPA, com e sem utilização de grade uniforme. Fonte: Autor.}
    \label{fig:tempo-comparativo}
\end{figure}

A Figura~\ref{fig:tempo-comparativo} apresenta o tempo médio de execução em função do número de objetos para as diferentes configurações avaliadas. Observa-se que os cenários sem utilização de grade uniforme apresentam crescimento acentuado, compatível com comportamento quadrático, enquanto as abordagens com filtragem demonstram crescimento significativamente mais lento.

Nota-se que as curvas correspondentes aos algoritmos SAT e GJK/EPA apresentam comportamento semelhante. Embora o tempo médio do SAT seja levemente inferior ao GJK/EPA, no geral a escolha da estratégia de filtragem possui maior impacto no desempenho global do sistema do que o algoritmo utilizado na fase de refinamento.

Observa-se também que, em cenários de alta densidade, o custo computacional cresce rapidamente na ausência de estruturas de filtragem, evidenciando a importância de técnicas de particionamento espacial para aplicações em tempo real.

\begin{figure}[htb]
    \centering
    \includegraphics[width=0.85\textwidth]{comparativo-tempo-execução-planckjs.png}
    \caption{Tempo médio de execução em função do número de objetos com a biblioteca \textit{Planck.js}. Fonte: Autor.}
    \label{fig:tempo-comparativo-planckjs}
\end{figure}

A Figura~\ref{fig:tempo-comparativo-planckjs} apresenta o tempo médio de execução em função do número de objetos para a biblioteca \textit{Planck.js}. Observa-se que o desempenho da \textit{Planck.js} apresenta comportamento consistente com o uso de uma estrutura eficiente na fase de filtragem, evidenciando crescimento controlado mesmo com o aumento do número de objetos. Esse resultado confirma a eficácia do uso de árvores hierárquicas dinâmicas na redução do número de pares de objetos avaliados.

De forma geral, os resultados indicam que a combinação de algoritmos eficientes na fase de \textit{narrow-phase} com estratégias adequadas de \textit{broad-phase} é essencial para garantir desempenho e escalabilidade em sistemas de detecção de colisão.